\documentclass[11pt]{amsart}

\usepackage[a4paper,margin=2.5cm]{geometry}
\usepackage{amsmath,amssymb,amsthm,mathtools}
\usepackage{enumitem}
\usepackage{setspace}
\usepackage[hidelinks]{hyperref}

\setstretch{1.2}

\newtheorem{theorem}{Theorem}[section]
\newtheorem{proposition}[theorem]{Proposition}
\newtheorem{lemma}[theorem]{Lemma}
\newtheorem{corollary}[theorem]{Corollary}
\theoremstyle{definition}
\newtheorem{definition}[theorem]{Definition}
\newtheorem{remark}[theorem]{Remark}

\newcommand{\R}{\mathbb R}
\newcommand{\Om}{\Omega}
\newcommand{\Fr}{\mathcal A}
\newcommand{\HH}{\mathcal H}
\newcommand{\dT}{\delta_T}
\newcommand{\rstar}{\rho_*}
\newcommand{\rdel}{\rho_\delta}
\newcommand{\wt}{w}
\newcommand{\eps}{\varepsilon}
\newcommand{\Hn}{\mathcal H^{n-1}}
\newcommand{\abs}[1]{\left|#1\right|}
\newcommand{\Per}{P}

\title[Moving-ball closure of the bad-density step]
{A moving-ball closure of the bad-density step\\
in the level-set route to sharp Saint--Venant stability}
\author{}
\date{May 2026}

\begin{document}

\begin{abstract}
For a bounded open set \(\Om\subset\R^n\) of volume \(\omega_n\) we prove the
bounded sharp Saint--Venant stability estimate
\(\Fr(\Om)^2\le C\,\dT(\Om)\), where \(\Fr\) is the Fraenkel asymmetry and
\(\dT\) the Saint--Venant deficit.  The argument is built around the unscaled
moving-ball distance \(q(\rho)=\inf_{z}|E_\rho\Delta B_\rho(z)|\) on the
torsion superlevel sets \(E_\rho\).  We first establish, on reduced boundaries
and for almost every level, the exact level-set deficit identity expressing
\(\dT\) as a weighted integral of the Cauchy--Schwarz defect \(D_H\) and the
isoperimetric defect \(D_I\), and derive from it a profile-gap identity and a
bad-density bound.  Nestedness gives a perimeter-free Lipschitz bound for
\(q\); we then prove a finite-perimeter one-sided affine first variation for
\(q\), a positive kinetic estimate, and an endpoint trace, which together close
the stability estimate.  The same-centre good-level package is obtained from
the Fusco--Julin strong quantitative isoperimetric inequality and an elementary
reduced-boundary radial trace lemma.  The manuscript is self-contained modulo
standard tools for sets of finite perimeter (the BV coarea formula, the
Gauss--Green theorem, and Reshetnyak continuity) and the cited isoperimetric
theorems.
\end{abstract}

\maketitle
\tableofcontents

\section{Introduction}\label{sec:intro}

Let \(\Om\subset\R^n\) (\(n\ge2\)) be a bounded open set with \(|\Om|=\omega_n\)
and let \(u=u_\Om\) be its torsion function.  Writing \(\dT(\Om)\) for the
Saint--Venant deficit and \(\Fr(\Om)\) for the Fraenkel asymmetry, the main
result is the bounded sharp stability estimate
\[
   \Fr(\Om)^2\le C(n,R,\rstar)\,\dT(\Om)
   \qquad(\Om\subset B_R),
\]
proved as Theorem~\ref{thm:boundedSV}.  The strategy works directly with the
torsion superlevel sets \(E_\rho=\{u>t(\rho)\}\) of volume \(\omega_n\rho^n\)
and the \emph{unscaled} moving-ball distance
\[
   q(\rho):=\inf_{z\in\R^n}|E_\rho\Delta B_\rho(z)| .
\]
Because the family \(\{E_\rho\}\) is nested, \(q\) is uniformly Lipschitz
(Lemma~\ref{lem:qLip}); this is what lets the sparse set of bad-density radii be
charged by Lebesgue measure rather than by the deficit.

The proof has three parts.  Section~\ref{sec:identity} proves, on reduced
boundaries and for a.e.\ level, the exact deficit identity
\(2\dT(\Om)=\int_0^{\|u\|_\infty}(D_H+D_I)/(n^2\omega_n^{2/n}\,m^{1-2/n})\,dt\),
where \(D_H\) is the Cauchy--Schwarz defect of \(|\nabla u|\) on
\(\partial^*E_t\) and \(D_I\) the squared-perimeter isoperimetric defect, and
Section~\ref{sec:profilegap} derives its profile-gap and bad-density
consequences.  Section~\ref{sec:samecentre} converts the Fusco--Julin strong
quantitative isoperimetric inequality into a same-centre control of the
homothetic trace.  Sections~\ref{sec:coarea}--\ref{sec:kinetic} prove a
finite-perimeter one-sided affine first variation for \(q\) and the resulting
positive kinetic estimate, and Section~\ref{sec:endpoint} closes the argument
through an endpoint trace.  The one substantive external theorem is the
Fusco--Julin quantitative isoperimetric inequality; the remaining ingredients
are classical tools for sets of finite perimeter and classical
rearrangement and isoperimetric inequalities, collected in
Appendix~\ref{app:bv} and the dependency certificate of
Appendix~\ref{app:constants}.

\section{Standing notation and scope}\label{sec:notation}

Fix \(n\ge2\), \(R\ge1\), and \(0<\rstar<1\).  Let
\(\Om\subset B_R\) be bounded and open with \(|\Om|=\omega_n\), and let
\(u=u_\Om\in H^1_0(\Om)\) be the torsion function,
\[
   -\Delta u=1\quad\hbox{in }\Om,\qquad u=0\quad\hbox{on }\partial\Om
   \quad(H^1_0\hbox{ sense}).
\]
We write the Saint--Venant deficit as
\[
   \dT(\Om):=E(\Om)-E(B_1)
   =-\frac12\int_\Om u+\frac12\int_{B_1}u_{B_1}\ge0.
\]
Let \(t(\rho)\) be the volume-radius parametrisation of the superlevels:
\[
   E_\rho:=\{u>t(\rho)\},\qquad |E_\rho|=\omega_n\rho^n,
   \qquad 0<\rho<1.
\]
Set
\[
   \wt(\rho):=-t'(\rho),\qquad d\mu(\rho)=\wt(\rho)\,d\rho.
\]
For a.e. \(\rho\), \(E_\rho\) has finite perimeter, \(\Hn(\partial^*E_\rho
\cap\{|\nabla u|=0\})=0\), and the De Giorgi outer normal satisfies
\[
   \nu_\rho=-\frac{\nabla u}{|\nabla u|}
   \quad\Hn\hbox{-a.e. on }\partial^*E_\rho
\]
(these level conventions are established in
Lemma~\ref{lem:id-noatoms}).  On those levels define
\[
   V_\rho:=\frac{\wt(\rho)}{|\nabla u|},\qquad
   H_{z,\rho}(x):=\frac{(x-z)\cdot\nu_\rho(x)}{\rho},\qquad
   \bar V_\rho:=\frac{\Per(B_\rho)}{\Per(E_\rho)}.
\]
The coarea relation of \S\ref{sec:identity}
(Lemma~\ref{lem:id-coarea}, with \(\wt(\rho)\,\alpha(t(\rho))=\Per(B_\rho)\))
gives
\begin{equation}\label{eq:V-mean}
   \int_{\partial^*E_\rho} V_\rho\,d\Hn=\Per(B_\rho),
   \qquad
   \bar V_\rho=\frac1{\Per(E_\rho)}
       \int_{\partial^*E_\rho}V_\rho\,d\Hn .
\end{equation}

The two level-set defects are
\[
   D_H(t(\rho))
   :=\left(\int_{\partial^*E_\rho}\frac1{|\nabla u|}\,d\Hn\right)
     \left(\int_{\partial^*E_\rho}|\nabla u|\,d\Hn\right)
      -\Per(E_\rho)^2,
\]
\[
   D_I(t(\rho)):=\Per(E_\rho)^2-\Per(B_\rho)^2 .
\]

\begin{remark}[Scope]\label{rem:scope}
The bounded Saint--Venant estimate is deduced from the exact level-set
identity (Theorem~\ref{thm:id-main}), proved in \S\ref{sec:identity}, and its
profile-gap and bad-density consequences (Proposition~\ref{prop:levelset}).
The same-centre good-level package (Theorem~\ref{inp:samecenter}) rests on the
theorem of Fusco--Julin, used in its exact same-centre form; the radial trace
conversion from normal oscillation to the homothetic trace is
Lemma~\ref{lem:oriented-radial-trace}.  Every other step, including the first
variation and the bad-density bookkeeping, is proved here.
\end{remark}

\section{Finite-perimeter level-set identities}\label{sec:identity}

All integrals in this section are taken on De~Giorgi reduced boundaries and all
statements hold for almost every level; no smoothness of \(\partial\Om\), no
graph parametrisation of the level sets, and no pointwise lower bound on
\(|\nabla u|\) are used.  We keep the normalisation \(|\Om|=\omega_n\), so the
ball of the same volume is \(B_1\); the general case follows by scaling.  Write
\[
   M:=\|u\|_\infty,\qquad L:=|\Om|=\omega_n,\qquad
   m(t):=|\{u>t\}|,
\]
\[
   \alpha(t):=\int_{\partial^*E_t}\frac1{|\nabla u|}\,d\Hn,\qquad
   \beta(t):=\int_{\partial^*E_t}|\nabla u|\,d\Hn,\qquad
   c_n:=n^2\omega_n^{2/n}.
\]
By the Talenti pointwise comparison \cite{Talenti1976},
\(u^*(s)\le u_{B_1}^*(s)=(1-(s/\omega_n)^{2/n})/(2n)\), so
\(M\le 1/(2n)<\infty\); here \(u^*\) is the decreasing rearrangement of \(u\).

\subsection*{Regularity and level conventions}

\begin{lemma}[Critical set and level conventions]\label{lem:id-noatoms}
Let \(u\) be extended by \(0\) to \(\R^n\).  Then:
\begin{enumerate}[label=\textup{(\alph*)},leftmargin=2.2em]
\item \(|\{u=0\}\cap\Om|=0\) and \(|\{\nabla u=0\}\cap\Om|=0\);
\item for every \(t\in(0,M)\), \(|\{u=t\}|=0\); hence \(m\) is continuous and
strictly decreasing on \((0,M)\);
\item for a.e.\ \(t\in(0,M)\): \(E_t\) has finite perimeter,
\(\partial^*E_t\subset\Om\) up to \(\Hn\)-null sets,
\(\nu_{E_t}=-\nabla u/|\nabla u|\) \(\Hn\)-a.e.\ on \(\partial^*E_t\), and
\(\Hn(\partial^*E_t\cap\{\nabla u=0\})=0\).
\end{enumerate}
\end{lemma}

\begin{proof}
\emph{(a)} Testing \(-\Delta u=1\) against \(u^-\in H^1_0(\Om)\) gives
\(\int|\nabla u^-|^2=-\int u^-\le0\), so \(u\ge0\) a.e.; the weak strong maximum
principle on each connected component yields \(u>0\) a.e.\ in \(\Om\)
\cite{GilbargTrudinger2001}, i.e.\ \(|\{u=0\}\cap\Om|=0\).  Interior elliptic
regularity gives \(u\in W^{2,p}_{\rm loc}(\Om)\) for all \(p<\infty\); applying
Stampacchia's theorem \cite{Stampacchia1965} to the weak derivatives
\(\partial_iu\) gives \(D^2u=0\) a.e.\ on \(\{\nabla u=0\}\), while
\(\Delta u=-1\) a.e.  Hence \(|\{\nabla u=0\}\cap\Om|=0\).

\emph{(b)} For \(t\in(0,M)\) the Sobolev chain rule applied to \((u-t)^+\) gives
\(\nabla u=0\) a.e.\ on \(\{u=t\}\); since \(\{u=t\}\subset\{u>0\}\subset\Om\) up
to a null set, part~(a) gives \(|\{u=t\}|=0\).  A jump of \(m\) at \(t\) would
have size \(|\{u=t\}|=0\), so \(m\) is continuous on \((0,M)\).  For
\(0<a<b<M\), some component of \(\Om\) has essential supremum of \(u\) exceeding
\(b\); on it \(\{a<u<b\}\) is a nonempty open set of positive measure, so
\(m(a)-m(b)=|\{a<u\le b\}|>0\).

\emph{(c)} By the BV coarea formula \cite[Thm.~13.1]{Maggi2012},
\cite[Thm.~3.40]{AFP2000}, for every Borel \(\varphi\ge0\),
\begin{equation}\label{eq:id-bvcoarea}
   \int_\Om|\nabla u|\,\varphi\,dx
   =\int_{-\infty}^{\infty}\Bigl(\int_{\partial^*E_t}\varphi\,d\Hn\Bigr)\,dt .
\end{equation}
Taking \(\varphi=\mathbf 1_{(0,M)}(u)\) (which equals \(1\) on
\(\partial^*E_t\) precisely for \(t\in(0,M)\)) gives
\(\int_0^M\Per(E_t)\,dt=\int_{\{0<u<M\}}|\nabla u|\,dx\le\int_\Om|\nabla u|
<\infty\), the last bound because \(|\nabla u|\in L^2(\Om)\subset L^1(\Om)\) on
the bounded set \(\Om\); hence \(\Per(E_t)<\infty\) for a.e.\ \(t\).  For \(t>0\) the precise representative of
\(u\) vanishes \(\Hn\)-a.e.\ on \(\R^n\setminus\Om\) (as \(u\in H^1_0(\Om)\)),
and for a.e.\ \(t\) one has \(\partial^*E_t\subset\{u=t\}\) up to \(\Hn\)-null
sets \cite{AFP2000,Maggi2012}; therefore \(\partial^*E_t\subset\Om\) up to
\(\Hn\)-null sets.  The orientation \(\nu_{E_t}=-\nabla u/|\nabla u|\) is the
standard one for superlevel sets of Sobolev functions \cite{AFP2000}.  Finally,
\eqref{eq:id-bvcoarea} with \(\varphi=\mathbf 1_{\{\nabla u=0\}}\) gives
\(\int_0^M\Hn(\partial^*E_t\cap\{\nabla u=0\})\,dt
=\int_{\{\nabla u=0\}}|\nabla u|\,dx=0\), so the last assertion holds for a.e.\
\(t\).
\end{proof}

\noindent
Statements ``for a.e.\ \(t\)'' below refer to the full-measure set of levels
furnished by Lemma~\ref{lem:id-noatoms}(c).

\subsection*{Coarea form of \texorpdfstring{\(-m'\)}{-m'} and the flux identity}

\begin{lemma}[Coarea form of \(-m'\)]\label{lem:id-coarea}
For a.e.\ \(t\in(0,M)\), \(-m'(t)=\alpha(t)\).  In particular \(m\) is locally
absolutely continuous on \((0,M)\) and \(\alpha\in L^1_{\rm loc}((0,M))\).
\end{lemma}

\begin{proof}
Fix \(0<a<b<M\) and apply \eqref{eq:id-bvcoarea} to
\(\varphi=|\nabla u|^{-1}\mathbf 1_{\{|\nabla u|>0\}}\mathbf 1_{(a,b)}(u)\).
By Lemma~\ref{lem:id-noatoms}(a) the integrand on the left equals
\(\mathbf 1_{\{a<u<b\}}\) up to a null set, so
\(m(a)-m(b)=\int_a^b\alpha(t)\,dt\).  Local absolute continuity of \(m\) and
\(-m'=\alpha\) a.e.\ follow.
\end{proof}

\begin{lemma}[Flux equals volume]\label{lem:id-flux}
For a.e.\ \(t\in(0,M)\), \(\beta(t)=\int_{\partial^*E_t}|\nabla u|\,d\Hn=m(t)\).
\end{lemma}

\begin{proof}
For \(\eps>0\) set \(\phi_\eps(s):=\min(1,\max(0,s/\eps))\) and
\(\eta_\eps:=\phi_\eps(u-t)\in H^1_0(\Om)\).  Testing \(-\Delta u=1\) against
\(\eta_\eps\),
\[
   \int_\Om\nabla u\cdot\nabla\eta_\eps\,dx=\int_\Om\eta_\eps\,dx .
\]
The Sobolev chain rule \cite[Thm.~3.99]{AFP2000} gives
\(\nabla\eta_\eps=\eps^{-1}\nabla u\,\mathbf 1_{\{t<u<t+\eps\}}\), so the left
side equals \(\eps^{-1}\int_{\{t<u<t+\eps\}}|\nabla u|^2\), which by
\eqref{eq:id-bvcoarea} (with \(\varphi=|\nabla u|\,\mathbf 1_{(t,t+\eps)}(u)\))
equals \(\eps^{-1}\int_t^{t+\eps}\beta(\tau)\,d\tau\).  As \(\eps\downarrow0\),
\(\eta_\eps\to\mathbf 1_{E_t}\) boundedly (Lemma~\ref{lem:id-noatoms}(b)), so the
right side tends to \(m(t)\).  Since \(\beta\in L^1((0,M))\) by
\eqref{eq:id-bvcoarea} and the Dirichlet identity
\(\int_0^M\beta=\int_\Om|\nabla u|^2=\int_\Om u<\infty\), at every Lebesgue
point of \(\beta\) the left side tends to \(\beta(t)\).  Hence
\(\beta(t)=m(t)\) for a.e.\ \(t\).
\end{proof}

\begin{remark}[Why a truncation, not a trace]\label{rem:id-flux}
Lemma~\ref{lem:id-flux} is the rigorous replacement, on an arbitrary
finite-perimeter level set, for the formal divergence computation
\(m(t)=-\int_{E_t}\Delta u=-\int_{\partial^*E_t}\nabla u\cdot\nu_{E_t}=\beta(t)\)
(using \(\nabla u\cdot\nu_{E_t}=-|\nabla u|\)); no boundary trace of
\(\nabla u\) on \(\partial^*E_t\) is required \cite[Ch.~17]{Maggi2012}.
\end{remark}

\subsection*{The rearranged primitive}

Let \(I(s):=\int_0^s u^*(\sigma)\,d\sigma\), \(s\in[0,L]\).

\begin{lemma}[Derivatives of the primitive]\label{lem:id-I}
\(I\) is absolutely continuous on \([0,L]\) with \(I'(s)=u^*(s)\) a.e.; on every
compact subinterval of \((0,L)\) the rearrangement \(u^*\) is absolutely
continuous and
\begin{equation}\label{eq:id-Ipp}
   I''(s)=(u^*)'(s)=\frac{1}{m'(u^*(s))}=-\frac{1}{\alpha(u^*(s))}
   \qquad\text{for a.e.\ }s\in(0,L).
\end{equation}
\end{lemma}

\begin{proof}
By the Talenti bound \(u^*\in L^\infty(0,L)\subset L^1(0,L)\); the Lebesgue
differentiation theorem gives \(I'=u^*\) a.e.  By Lemma~\ref{lem:id-coarea},
\(m\) is locally absolutely continuous with \(m'=-\alpha\) a.e., and \(m\) is
continuous and strictly decreasing on \((0,M)\) by
Lemma~\ref{lem:id-noatoms}(b), so \(m(u^*(s))=s\).  The absolutely continuous
inverse theorem for monotone functions \cite[Thm.~0.4, \S1]{BBC1975} gives that
\(u^*\) is absolutely continuous on compact subintervals of \((0,L)\) with
\((u^*)'(s)=1/m'(u^*(s))\) wherever \(m'(u^*(s))\) exists and is nonzero.  For
a.e.\ \(t\in(0,M)\) one has \(0<|E_t|<|\Om|\), hence \(\Per(E_t)>0\) and
\(\alpha(t)>0\); thus \(\{m'=0\}=\{\alpha=0\}\) is null and carries no image
measure (\(|m(\{m'=0\})|\le\int_{\{m'=0\}}|m'|=0\)).  Since \(m'=-\alpha\),
\eqref{eq:id-Ipp} follows; absolute continuity of \(u^*\) on compact
subintervals identifies the distributional and a.e.\ derivatives of
\(I'=u^*\).
\end{proof}

\begin{remark}[The corrected second derivative]\label{rem:id-alpha}
The factor in \eqref{eq:id-Ipp} is \(\alpha=\int_{\partial^*E_t}|\nabla u|^{-1}\),
not \(\beta=\int_{\partial^*E_t}|\nabla u|\).  Since \(\beta(t)=m(t)\)
(Lemma~\ref{lem:id-flux}), replacing \(\alpha\) by \(\beta\) would give
\(I''=-1/m\), which holds only when \(|\nabla u|\) is constant on
\(\partial^*E_t\), i.e.\ on the ball; the discrepancy is exactly the
Cauchy--Schwarz defect \(D_H\) introduced next, so the substitution would
silently delete \(D_H\) from the identity.
\end{remark}

\subsection*{The exact deficit identity}

For a.e.\ \(t\in(0,M)\) the two level-set defects of the standing notation read,
as functions of the level,
\[
   D_H(t)=\alpha(t)\beta(t)-\Per(E_t)^2,\qquad
   D_I(t)=\Per(E_t)^2-c_n\,m(t)^{2-2/n},
\]
consistent with \S\ref{sec:notation} under \(t=t(\rho)\) (so
\(\Per(B_\rho)^2=c_n m(t)^{2-2/n}\)).  By Cauchy--Schwarz on
\(\partial^*E_t\) applied to \((|\nabla u|^{1/2},|\nabla u|^{-1/2})\),
\(\Per(E_t)^2\le\alpha(t)\beta(t)\), so \(D_H(t)\ge0\), with equality iff
\(|\nabla u|\) is \(\Hn\)-a.e.\ constant on \(\partial^*E_t\).  By the De~Giorgi
isoperimetric inequality \cite{DeGiorgi1958,Maggi2012},
\(\Per(E_t)\ge n\omega_n^{1/n}m(t)^{1-1/n}\), so \(D_I(t)\ge0\), with equality
iff \(E_t\) is equivalent to a ball.

\begin{lemma}[Pointwise convexity formula]\label{lem:id-Gpp}
Set \(G(s):=I(s)+\dfrac{s^{1+2/n}}{(4+2n)\,\omega_n^{2/n}}\).  Then for a.e.\
\(s\in(0,L)\), with \(t=u^*(s)\),
\[
   G''(s)=\frac{1}{c_n\,s^{1-2/n}}-\frac{1}{\alpha(t)} .
\]
\end{lemma}

\begin{proof}
Direct differentiation gives
\(\frac{d^2}{ds^2}\frac{s^{1+2/n}}{(4+2n)\omega_n^{2/n}}
=\frac{1}{c_n s^{1-2/n}}\); add \eqref{eq:id-Ipp}.
\end{proof}

\begin{lemma}[Convexity identity]\label{lem:id-cov}
\(\displaystyle\int_0^L s\,G''(s)\,ds
   =\int_0^M\frac{D_H(t)+D_I(t)}{c_n\,m(t)^{1-2/n}}\,dt.\)
\end{lemma}

\begin{proof}
Substitute \(s=m(t)\), \(ds=-\alpha(t)\,dt\), in
\(\int_0^L s\,G''(s)\,ds\) (Lemma~\ref{lem:id-noatoms}(b) removes jump terms,
and \(\{\alpha=0\}\) carries no image measure):
\[
   \int_0^L s\,G''\,ds
   =\int_0^M m(t)\Bigl(\frac{1}{c_n m(t)^{1-2/n}}-\frac1{\alpha(t)}\Bigr)
      \alpha(t)\,dt
   =\int_0^M\Bigl(\frac{m(t)\alpha(t)}{c_n m(t)^{1-2/n}}-m(t)\Bigr)dt .
\]
Use \(m(t)=\beta(t)\) (Lemma~\ref{lem:id-flux}) and add and subtract
\(\Per(E_t)^2/(c_n m(t)^{1-2/n})\):
\[
   \frac{m\alpha}{c_n m^{1-2/n}}-m
   =\frac{\alpha\beta-c_n m^{2-2/n}}{c_n m^{1-2/n}}
   =\frac{D_H+D_I}{c_n m^{1-2/n}} . \qedhere
\]
\end{proof}

\begin{lemma}[Endpoint convexity-gap formula]\label{lem:id-gap}
With \(\Gamma(\Om):=G'_-(L)-G(L)/L\), one has
\(\Gamma(\Om)=\frac1L\int_0^L s\,G''(s)\,ds\).
\end{lemma}

\begin{proof}
\(G\in W^{2,1}_{\rm loc}((0,L])\) and \(G(0)=I(0)=0\).  Integration by parts on
\((\eps,L)\) gives \(\int_\eps^L sG''=L\,G'_-(L)-\eps G'(\eps)-G(L)+G(\eps)\).
By Lemma~\ref{lem:id-I},
\(G'(s)=u^*(s)+s^{2/n}/(2n\omega_n^{2/n})\) is bounded, so
\(\eps G'(\eps)\to0\) and \(G(\eps)\to0\); divide by \(L\).
\end{proof}

\begin{lemma}[Endpoint value]\label{lem:id-endpoint}
\(\Gamma(\Om)=\dfrac{2}{|\Om|}\bigl(E(\Om)-E(B_1)\bigr)=\dfrac{2}{\omega_n}\dT(\Om).\)
\end{lemma}

\begin{proof}
At \(s=L\), \(u^*(s)\to0\) as \(s\uparrow L\) (Lemma~\ref{lem:id-noatoms} gives
\(m(t)\uparrow L\) as \(t\downarrow0\)), so \(G'_-(L)=L^{2/n}/(2n\omega_n^{2/n})\)
and \(I(L)=\int_\Om u\,dx\).  With \(L=\omega_n\), using
\(\frac{1}{2n}-\frac1{4+2n}=\frac{1}{2n}-\frac1{2(n+2)}=\frac1{n(n+2)}\),
\[
   \Gamma(\Om)
   =\frac{L^{2/n}}{2n\omega_n^{2/n}}-\frac{L^{2/n}}{(4+2n)\omega_n^{2/n}}
    -\frac1L\int_\Om u
   =\frac{1}{n(n+2)}-\frac1{\omega_n}\int_\Om u .
\]
For the unit-volume ball \(B_1\), \(u_{B_1}(x)=(1-|x|^2)/(2n)\) and
\[
   \int_{B_1}u_{B_1}
   =\frac1{2n}\,n\omega_n\int_0^1(1-r^2)r^{n-1}\,dr
   =\frac{\omega_n}{2}\Bigl(\frac1n-\frac1{n+2}\Bigr)
   =\frac{\omega_n}{n(n+2)},
\]
so \(\frac1{\omega_n}\int_{B_1}u_{B_1}=\frac1{n(n+2)}\) and \(\Gamma(B_1)=0\)
(the P\'olya--Szeg\H{o} identification \cite{PolyaSzego1951,Bandle1980}).
Subtracting, \(\Gamma(\Om)=\frac1{\omega_n}(\int_{B_1}u_{B_1}-\int_\Om u)
=\frac{2}{\omega_n}(E(\Om)-E(B_1))\), since \(E(\cdot)=-\tfrac12\int(\cdot)\).
\end{proof}

\begin{theorem}[Exact deficit identity on reduced-boundary level sets]
\label{thm:id-main}
Let \(\Om\subset\R^n\) be open with \(|\Om|=\omega_n\) and \(u\in H^1_0(\Om)\)
the torsion function.  Then \eqref{eq:id-bvcoarea} and
Lemmas~\ref{lem:id-coarea}, \ref{lem:id-flux}, \ref{lem:id-I} hold, and
\begin{equation}\label{eq:id-main}
   \boxed{\;
   2\,\dT(\Om)
   =\int_0^{\|u\|_\infty}
      \frac{D_H(t)+D_I(t)}{n^2\omega_n^{2/n}\,m(t)^{1-2/n}}\,dt ,\;}
\end{equation}
with \(D_H(t),D_I(t)\ge0\) for a.e.\ \(t\).  By scaling, the corresponding
identity \(\tfrac{2}{|\Om|}(E(\Om)-E(B))=\tfrac1{|\Om|}\int_0^{\|u\|_\infty}
(D_H+D_I)/(c_n m^{1-2/n})\,dt\) holds for every open \(\Om\) of finite measure.
\end{theorem}

\begin{proof}
Combine Lemmas~\ref{lem:id-cov}, \ref{lem:id-gap}, \ref{lem:id-endpoint}; the
scaling statement follows since both sides are \(0\)-homogeneous under
\(\Om\mapsto\lambda\Om\) after the displayed normalisation.
\end{proof}

\section{Profile gap and bad-density estimate}\label{sec:profilegap}

\begin{proposition}[Level-set budget and profile-gap bad set]
\label{prop:levelset}
There are computable constants
\[
   C_L=C_L(n,\rstar),\qquad C_\tau=C_\tau(n,\rstar),
   \qquad \kappa=\kappa(\rstar)>0,
\]
such that, with
\[
   \rdel:=1-\kappa\sqrt{\dT(\Om)}
\]
and whenever \(\rdel\ge(1+\rstar)/2\), one has
\begin{equation}\label{eq:budget}
   \int_{\rstar}^{\rdel}
      \bigl(D_H(t(\rho))+D_I(t(\rho))\bigr)\,d\mu(\rho)
   \le C_L\,\dT(\Om),
\end{equation}
\[
   0\le \wt(\rho)\le\frac1n\quad\hbox{for a.e. }\rho,
\]
and, with
\[
   \tau_0:=\frac{\rstar}{2n},\qquad
   B_\tau:=\{\rho\in[\rstar,\rdel]:\wt(\rho)<\tau_0\},\qquad
   G_\tau:=[\rstar,\rdel]\setminus B_\tau,
\]
one has
\begin{equation}\label{eq:badtau}
   |B_\tau|\le C_\tau\,\dT(\Om).
\end{equation}
\end{proposition}

\begin{proof}
Recall \(m(t)=|\{u>t\}|\) and
\(\alpha(t)=\int_{\partial^*\{u>t\}}|\nabla u|^{-1}\,d\Hn\) from
\S\ref{sec:identity}.  Theorem~\ref{thm:id-main} gives the exact deficit
identity
\begin{equation}\label{eq:levelset-identity}
   2\,\dT(\Om)
   =
   \int_0^{\|u\|_\infty}
      \frac{D_H(t)+D_I(t)}
      {n^2\omega_n^{2/n}m(t)^{1-2/n}}\,dt ,
\end{equation}
all integrals taken on the reduced boundaries \(\partial^*\{u>t\}\) and holding
for a.e.\ level.

Changing variables \(t=t(\rho)\), \(m(t(\rho))=\omega_n\rho^n\), and
\(dt=d\mu(\rho)\) in decreasing-radius orientation gives
\[
   2\dT(\Om)
   =
   \int_0^1
      \frac{D_H(t(\rho))+D_I(t(\rho))}
      {n^2\omega_n\,\rho^{n-2}}\,d\mu(\rho).
\]
Since \(n\ge2\) and \(0<\rho\le1\),
\[
   \int_{\rstar}^{\rdel}
      \bigl(D_H(t(\rho))+D_I(t(\rho))\bigr)\,d\mu(\rho)
   \le 2n^2\omega_n\,\dT(\Om),
\]
which proves \eqref{eq:budget}, for instance with
\(C_L=2n^2\omega_n\).

It remains to prove the density bounds.  Write
\[
   a(\rho):=\wt(\rho)=-t'(\rho),\qquad a_B(\rho):=\frac{\rho}{n}.
\]
We next derive from \eqref{eq:levelset-identity} the profile
derivative gap identity
\begin{equation}\label{eq:profile-gap-rho}
   0\le a(\rho)\le a_B(\rho)
   \quad\hbox{for a.e. }\rho\in(0,1),
   \qquad
   \int_0^1 \rho^n\bigl(a_B(\rho)-a(\rho)\bigr)\,d\rho
   =
   \frac{2}{\omega_n}\dT(\Om).
\end{equation}
For completeness, the algebra is as follows.  Coarea gives
\[
   a(\rho)=\frac{n\omega_n\rho^{n-1}}{\alpha(t(\rho))}.
\]
For the ball profile the corresponding coarea coefficient is
\[
   \alpha_B(\rho)=n^2\omega_n\rho^{n-2},
   \qquad
   \frac{n\omega_n\rho^{n-1}}{\alpha_B(\rho)}=\frac{\rho}{n}=a_B(\rho).
\]
Since
\[
   D_H(t(\rho))+D_I(t(\rho))
   =\omega_n\rho^n\bigl(\alpha(t(\rho))-\alpha_B(\rho)\bigr),
\]
the nonnegativity of \(D_H+D_I\) gives \(0\le a\le a_B\), and inserting
the same formula in \eqref{eq:levelset-identity} gives the
integral identity in \eqref{eq:profile-gap-rho}.

Thus \(0\le\wt(\rho)\le\rho/n\le1/n\).  On
\[
   B_\tau=\left\{\rho\in[\rstar,\rdel]:
      a(\rho)<\frac{\rstar}{2n}\right\}
\]
one has
\[
   a_B(\rho)-a(\rho)\ge \frac{\rho}{n}-\frac{\rstar}{2n}
      \ge\frac{\rstar}{2n}.
\]
Using \(\rho^n\ge\rstar^n\) on \([\rstar,\rdel]\) and
\eqref{eq:profile-gap-rho},
\[
   \frac{\rstar^{n+1}}{2n}|B_\tau|
   \le
   \int_{B_\tau}\rho^n(a_B-a)\,d\rho
   \le \frac{2}{\omega_n}\dT(\Om).
\]
Therefore \eqref{eq:badtau} holds with
\[
   C_\tau=\frac{4n}{\omega_n\rstar^{n+1}}.
\]
Finally, choose for example
\(\kappa=(1-\rstar)/4\).  Then \(\rdel\ge(1+\rstar)/2\) whenever
\(\dT(\Om)\le1\), and in the large-deficit part of the final theorem this
smallness is absorbed into the constant.
\end{proof}

\section{Quantitative isoperimetry and same-centre trace control}
\label{sec:samecentre}

\begin{theorem}[Fusco--Julin strong quantitative isoperimetry]
\label{thm:FJ-strong}
For every \(n\ge2\) there is a constant \(C_{\rm FJ}(n)\) with the
following property.  Let \(E\subset\R^n\) be a set of finite perimeter and
let \(r_E>0\) be defined by \(|E|=\omega_n r_E^n\).  Then there is a
centre \(z_E\in\R^n\) such that
\begin{equation}\label{eq:FJ-exact}
   \left[
      \frac{|E\Delta B_{r_E}(z_E)|}{r_E^n}
      +
      \left(
         \frac1{r_E^{\,n-1}}
         \int_{\partial^*E}
            \left|\nu_E(x)-\frac{x-z_E}{|x-z_E|}\right|^2
         d\Hn(x)
      \right)^{1/2}
   \right]^2
   \le
   C_{\rm FJ}(n)\,
   \frac{\Per(E)-\Per(B_{r_E})}{r_E^{\,n-1}} .
\end{equation}
\end{theorem}

\begin{proof}
This is exactly \cite[Theorem~1.1]{FJ2014}, with the Fusco--Julin
asymmetry
\[
   A(E)=\min_z\left\{
      \frac{|E\Delta B_{r_E}(z)|}{r_E^n}
      +
      \left(
         \frac1{r_E^{\,n-1}}
         \int_{\partial^*E}
            \left|\nu_E(x)-\nu_{B_{r_E}(z)}
               (\pi_{z,r_E}(x))\right|^2d\Hn(x)
      \right)^{1/2}
   \right\}
\]
and their deficit
\((\Per(E)-\Per(B_{r_E}))/r_E^{\,n-1}\).  Since
\(\nu_{B_{r_E}(z)}(\pi_{z,r_E}(x))=(x-z)/|x-z|\), their minimising centre
is one centre for both the symmetric-difference term and the boundary-normal
oscillation term.
\end{proof}

\begin{corollary}[Unnormalised same-centre consequences]
\label{cor:FJ-unnormalised}
Fix \(0<\rstar<1\).  Let \(E\subset\R^n\) be a set of finite perimeter with
\(|E|=\omega_n\rho^n\), where \(\rho\in[\rstar,1]\).  Then there is a
centre \(z\in\R^n\) such that
\begin{equation}\label{eq:FJ-volume}
   |E\Delta B_\rho(z)|^2
   \le C(n)\bigl(\Per(E)-\Per(B_\rho)\bigr),
\end{equation}
and
\begin{equation}\label{eq:FJ-normal}
   \int_{\partial^*E}
      \left|\nu_E(x)-\frac{x-z}{|x-z|}\right|^2d\Hn(x)
   \le C(n)\bigl(\Per(E)-\Per(B_\rho)\bigr).
\end{equation}
Consequently, since \(\rho\ge\rstar\),
\begin{equation}\label{eq:FJ-DI-version}
   |E\Delta B_\rho(z)|^2
   +
   \int_{\partial^*E}
      \left|\nu_E(x)-\frac{x-z}{|x-z|}\right|^2d\Hn(x)
   \le C(n,\rstar)
      \bigl(\Per(E)^2-\Per(B_\rho)^2\bigr).
\end{equation}
\end{corollary}

\begin{proof}
Apply Theorem~\ref{thm:FJ-strong} with \(r_E=\rho\).  Each summand in the
left side of \eqref{eq:FJ-exact} is nonnegative, so
\[
   \frac{|E\Delta B_\rho(z)|^2}{\rho^{2n}}
   +
   \frac1{\rho^{n-1}}
   \int_{\partial^*E}
      \left|\nu_E-\frac{x-z}{|x-z|}\right|^2d\Hn
   \le
   C(n)\frac{\Per(E)-\Per(B_\rho)}{\rho^{n-1}} .
\]
Since \(\rho\le1\), this gives \eqref{eq:FJ-volume} and
\eqref{eq:FJ-normal}.  Finally,
\[
   \Per(E)-\Per(B_\rho)
   =
   \frac{\Per(E)^2-\Per(B_\rho)^2}{\Per(E)+\Per(B_\rho)}
   \le
   \frac{\Per(E)^2-\Per(B_\rho)^2}{\Per(B_\rho)}
   \le C(n,\rstar)\bigl(\Per(E)^2-\Per(B_\rho)^2\bigr).
\]
\end{proof}

\begin{lemma}[Oriented radial trace]\label{lem:oriented-radial-trace}
Fix \(R\ge1\) and \(0<\rstar<1\).  Let \(E\subset B_R\) be a set of finite
perimeter with \(|E|=\omega_n\rho^n\), where \(\rho\in[\rstar,1]\), and
let \(z\in\R^n\) satisfy \(|z|\le R+1\).  Set
\[
   e_z(x):=\frac{x-z}{|x-z|},\qquad
   W_z(x):=\left|\frac{|x-z|}{\rho}-1\right|.
\]
Then
\begin{equation}\label{eq:oriented-radial-trace}
   \int_{\partial^*E} W_z\,d\Hn
   \le
   \frac{n}{\rstar}|E\Delta B_\rho(z)|
   +
   C(R,\rstar)
   \int_{\partial^*E}|\nu_E-e_z|^2\,d\Hn .
\end{equation}
\end{lemma}

\begin{proof}
Translate the notation temporarily so that \(z=0\), and write
\(B=B_\rho\), \(r=|x|\), \(e=x/r\).  The point \(x=0\) is
\(\Hn\)-negligible on \(\partial^*E\), so the value of \(e\) there is
irrelevant.  Define the bounded radial field
\[
   X(x):=\left|\frac{r}{\rho}-1\right|e .
\]
For \(0<r<\rho\),
\[
   \operatorname{div}X=\frac{n-1}{r}-\frac n\rho,
\]
while for \(r>\rho\),
\[
   \operatorname{div}X=\frac n\rho-\frac{n-1}{r}.
\]
The field is singular only at the centre and has a Lipschitz radial
coefficient away from the centre.  To justify Gauss--Green, replace \(X\)
by \(\chi_\eps(r)X\), where \(\chi_\eps=0\) on \(B_\eps\) and
\(\chi_\eps=1\) outside \(B_{2\eps}\), smooth the radial profile, apply the
classical finite-perimeter Gauss--Green formula \cite{Maggi2012}, and let the
smoothing and then \(\eps\downarrow0\).  The extra interior term
\(\int_E\nabla\chi_\eps\cdot X\,dx\) is supported on the annulus
\(\{\eps<|x|<2\eps\}\), where \(|\nabla\chi_\eps|\lesssim\eps^{-1}\) and
\(|X|\) is bounded; since that annulus has volume \(O(\eps^n)\), this term is
\(O(\eps^{n-1})\) and vanishes as \(\eps\downarrow0\).  The boundary error
tends to zero because
\(\Hn(\partial^*E\cap B_{2\eps})\to \Hn(\partial^*E\cap\{0\})=0\), and the
volume terms converge by dominated convergence, since \(r^{-1}\) is locally
integrable in \(\R^n\) for \(n\ge2\).  Thus
\[
   \int_{\partial^*E}X\cdot\nu_E\,d\Hn
   =
   \int_E\operatorname{div}X\,dx .
\]
The same identity for \(B\) gives zero, because \(X=0\) on \(\partial B\).
Subtracting it yields
\[
   \int_{\partial^*E}X\cdot\nu_E\,d\Hn
   =
   \int_{E\setminus B}\operatorname{div}X\,dx
   -
   \int_{B\setminus E}\operatorname{div}X\,dx .
\]
On \(E\setminus B\), \(\operatorname{div}X\le n/\rho\).  On
\(B\setminus E\), \(-\operatorname{div}X\le n/\rho\).  Hence
\begin{equation}\label{eq:radial-oriented-piece}
   \int_{\partial^*E}X\cdot\nu_E\,d\Hn
   \le \frac n\rho |E\Delta B|.
\end{equation}

On \(\partial^*E\),
\[
   W_z = X\cdot\nu_E + W_z(1-e_z\cdot\nu_E)
       = X\cdot\nu_E + \frac{W_z}{2}|\nu_E-e_z|^2 .
\]
Since \(E\subset B_R\) and \(|z|\le R+1\), every reduced-boundary point
satisfies \(|x-z|\le2R+1\), and therefore
\[
   W_z(x)\le 1+\frac{2R+1}{\rstar}.
\]
Combining this bound with \eqref{eq:radial-oriented-piece} and using
\(\rho\ge\rstar\) proves \eqref{eq:oriented-radial-trace}.
\end{proof}

\begin{theorem}[Same-centre good-level geometric package]\label{inp:samecenter}
There exist constants
\[
   \eta_I=\eta_I(n,R,\rstar)>0,\qquad
   C_{\rm sc}=C_{\rm sc}(n,R,\rstar),
\]
such that the following holds.  Define
\[
   G_I:=\{\rho\in[\rstar,\rdel]:D_I(t(\rho))\le\eta_I\}.
\]
For a.e. \(\rho\in G_I\) there is a centre \(z_\rho^{\rm sc}\) such that
\begin{equation}\label{eq:sc-input}
   |E_\rho\Delta B_\rho(z_\rho^{\rm sc})|^2
   +
   \left(
      \int_{\partial^*E_\rho}
         |H_{z_\rho^{\rm sc},\rho}-\bar V_\rho|\,d\Hn
   \right)^2
   \le C_{\rm sc}\,D_I(t(\rho)).
\end{equation}
\end{theorem}

\begin{proof}
Let \(\rho\in G_I\) be a finite-perimeter level.  Apply
Corollary~\ref{cor:FJ-unnormalised} to \(E_\rho\), and let
\(z=z_\rho^{\rm sc}\) be the centre obtained there.  Since
\[
   D_I(t(\rho))=\Per(E_\rho)^2-\Per(B_\rho)^2,
\]
\eqref{eq:FJ-DI-version} gives
\begin{equation}\label{eq:sc-proof-FJ}
   |E_\rho\Delta B_\rho(z)|^2
   +
   \int_{\partial^*E_\rho}|\nu_\rho-e_z|^2\,d\Hn
   \le C(n,\rstar)D_I(t(\rho)).
\end{equation}
Choose \(\eta_I\) so small that the first term in
\eqref{eq:sc-proof-FJ} implies
\[
   |E_\rho\Delta B_\rho(z)|<|B_\rho|.
\]
Then \(E_\rho\cap B_\rho(z)\) has positive measure.  Since
\(E_\rho\subset B_R\), this forces \(B_\rho(z)\cap B_R\ne\emptyset\), and
therefore
\begin{equation}\label{eq:sc-centre-localised}
   |z|\le R+\rho\le R+1 .
\end{equation}

The radial trace lemma gives
\[
   \int_{\partial^*E_\rho}
      \left|\frac{|x-z|}{\rho}-1\right|\,d\Hn
   \le
   C(n,R,\rstar)
   \bigl(D_I(t(\rho))^{1/2}+D_I(t(\rho))\bigr).
\]
Reducing \(\eta_I\le1\), this is bounded by
\begin{equation}\label{eq:radial-trace-DI}
   C(n,R,\rstar)D_I(t(\rho))^{1/2}.
\end{equation}
For \(H_{z,\rho}=(|x-z|/\rho)(e_z\cdot\nu_\rho)\),
\[
   |H_{z,\rho}-1|
   \le
   \left|\frac{|x-z|}{\rho}-1\right|
   +
   |1-e_z\cdot\nu_\rho|
   =
   \left|\frac{|x-z|}{\rho}-1\right|
   +
   \frac12|\nu_\rho-e_z|^2 .
\]
Together with \eqref{eq:sc-proof-FJ} and \eqref{eq:radial-trace-DI},
this gives
\begin{equation}\label{eq:Hminus1-DI}
   \int_{\partial^*E_\rho}|H_{z,\rho}-1|\,d\Hn
   \le C(n,R,\rstar)D_I(t(\rho))^{1/2}.
\end{equation}
Finally,
\[
   \bar V_\rho=\frac{\Per(B_\rho)}{\Per(E_\rho)},\qquad
   \Per(E_\rho)|1-\bar V_\rho|
   =
   \Per(E_\rho)-\Per(B_\rho)
   \le C(n,\rstar)D_I(t(\rho)).
\]
Combining this with \eqref{eq:Hminus1-DI} and again using
\(\eta_I\le1\),
\[
   \int_{\partial^*E_\rho}|H_{z,\rho}-\bar V_\rho|\,d\Hn
   \le C(n,R,\rstar)D_I(t(\rho))^{1/2}.
\]
Squaring and adding the first estimate in \eqref{eq:sc-proof-FJ} proves
\eqref{eq:sc-input}.
\end{proof}

\begin{lemma}[Velocity defect controlled by \(D_H\)]\label{lem:velocity}
For a.e. \(\rho\in[\rstar,\rdel]\),
\begin{equation}\label{eq:velocity-defect}
   \left(
      \int_{\partial^*E_\rho}|V_\rho-\bar V_\rho|\,d\Hn
   \right)^2
   \le D_H(t(\rho)).
\end{equation}
\end{lemma}

\begin{proof}
Write \(P=\Per(E_\rho)\), \(P^*=\Per(B_\rho)\),
\(\alpha=\int_{\partial^*E_\rho}|\nabla u|^{-1}\,d\Hn\), and
\(\beta=\int_{\partial^*E_\rho}|\nabla u|\,d\Hn\).  The coarea
parametrisation gives \(\wt\alpha=P^*\), and the flux identity gives
\(\beta=|E_\rho|=\omega_n\rho^n\).  With
\(f:=V_\rho=\wt/|\nabla u|\) and \(\bar f:=\bar V_\rho=P^*/P\),
\[
   \int_{\partial^*E_\rho} f\,d\Hn=P^*.
\]
Moreover
\[
\begin{aligned}
   \int_{\partial^*E_\rho}\frac{(f-\bar f)^2}{f}\,d\Hn
   &=\int f\,d\Hn-2\bar f P+\bar f^2\int \frac1f\,d\Hn  \\
   &=-P^*+\frac{(P^*)^2}{P^2}\frac{\beta}{\wt}
    =-P^*+\frac{P^*}{P^2}\alpha\beta                 \\
   &=\frac{P^*}{P^2}\bigl(\alpha\beta-P^2\bigr)
    =\frac{P^*}{P^2}D_H(t(\rho)).
\end{aligned}
\]
By Cauchy--Schwarz,
\[
   \left(\int |f-\bar f|\,d\Hn\right)^2
   \le
   \left(\int\frac{(f-\bar f)^2}{f}\,d\Hn\right)
   \left(\int f\,d\Hn\right)
   =
   \left(\frac{P^*}{P}\right)^2D_H(t(\rho)).
\]
The isoperimetric inequality gives \(P^*\le P\), hence
\eqref{eq:velocity-defect}.
\end{proof}

\begin{lemma}[Superlevel asymmetry transfer]\label{lem:superlevel-transfer}
Let \(E\subset\Om\), \(|\Om|=\omega_n\), and \(|E|=\omega_n\rho^n\) with
\(\rho\in(0,1]\).  Then
\[
   \Fr(\Om)\le \Fr(E)+\frac{2}{\omega_n}|\Om\setminus E|.
\]
Here \(\Fr(E)\) is the Fraenkel asymmetry of \(E\) relative to balls of
volume \(|E|\).
\end{lemma}

\begin{proof}
Let \(s=|E|\), \(L=|\Om|=\omega_n\), and \(\eta=L-s=|\Om\setminus E|\).
Choose \(z\) such that, for an arbitrary \(\eps>0\),
\[
   |E\Delta B_\rho(z)|\le s(\Fr(E)+\eps).
\]
Since \(B_\rho(z)\subset B_1(z)\),
\[
\begin{aligned}
   |\Om\Delta B_1(z)|
   &\le |\Om\Delta E|+|E\Delta B_\rho(z)|+|B_\rho(z)\Delta B_1(z)|\\
   &=\eta+|E\Delta B_\rho(z)|+\eta .
\end{aligned}
\]
Divide by \(L\), use \(s\le L\), and let \(\eps\downarrow0\).
\end{proof}

\section{Coarea differentiation and shell-admissible radii}\label{sec:coarea}

The first variation is proved on a fixed full-measure set of good-density
radii.  The point of the definition below is that all coarea differentiations
are registered in advance by a countable family independent of the radius
being tested.

\begin{lemma}[Two-sided variable-thickness coarea differentiation]
\label{lem:varcoarea}
Let \(\xi\le\zeta\) be bounded Borel functions on \(\R^n\).  Then for
Lebesgue-a.e. positive level \(s\in(0,\|u\|_\infty)\), with exceptional null
set depending only on the pair \((\xi,\zeta)\),
\begin{equation}\label{eq:varcoarea}
   \lim_{h\downarrow0}\frac1h
   \bigl|\{x:h\,\xi(x)<u(x)-s<h\,\zeta(x)\}\bigr|
   =
   \int_{\partial^*\{u>s\}}\frac{\zeta-\xi}{|\nabla u|}\,d\Hn .
\end{equation}
The statement is used only at levels for which the right side is finite and
\(\Hn(\partial^*\{u>s\}\cap\{|\nabla u|=0\})=0\).
\end{lemma}

\begin{proof}
By Lemma~\ref{lem:id-noatoms}(a), \(|\{x\in\Om:|\nabla u(x)|=0\}|=0\), so the
critical set may be discarded in the volume computations below.

Suppose first that \(\xi=\sum_j p_j{\bf1}_{A_j}\) and
\(\zeta=\sum_j q_j{\bf1}_{A_j}\) are simple over a common disjoint Borel
partition, with \(p_j\le q_j\).  By the BV coarea formula applied to
\({\bf1}_{A_j}|\nabla u|^{-1}{\bf1}_{\{|\nabla u|>0\}}\),
\[
   \bigl|\{hp_j<u-s<hq_j\}\cap A_j\bigr|
   =
   \int_{s+hp_j}^{s+hq_j}
      \int_{\partial^*\{u>\tau\}}\frac{{\bf1}_{A_j}}{|\nabla u|}
      \,d\Hn\,d\tau .
\]
The interval of integration shrinks to \(s\) and has length
\(h(q_j-p_j)\); hence at a Lebesgue point of
\(\tau\mapsto\int_{\partial^*\{u>\tau\}}{\bf1}_{A_j}|\nabla u|^{-1}\,d\Hn\),
division by \(h\) followed by \(h\downarrow0\) gives
\((q_j-p_j)\int_{\partial^*\{u>s\}}{\bf1}_{A_j}|\nabla u|^{-1}\,d\Hn\).
Summing over \(j\) proves \eqref{eq:varcoarea} for simple pairs.

For bounded \(\xi\le\zeta\), fix \(\eta>0\) and choose simple pairs that
sandwich the slab.  Outer: simple \(\sigma\le\xi\le\zeta\le\sigma'\) with
\(\sigma'-\sigma\le(\zeta-\xi)+2\eta\); then
\(\{h\xi<u-s<h\zeta\}\subseteq\{h\sigma<u-s<h\sigma'\}\), so the
\(\limsup\) of the left side of \eqref{eq:varcoarea} is at most
\[
   \int_{\partial^*\{u>s\}}\frac{\sigma'-\sigma}{|\nabla u|}\,d\Hn
   \le
   \int_{\partial^*\{u>s\}}\frac{\zeta-\xi}{|\nabla u|}\,d\Hn
   +2\eta\int_{\partial^*\{u>s\}}\frac1{|\nabla u|}\,d\Hn .
\]
Inner: simple \(\sigma_*\ge\xi\),
\(\sigma_*'\le\zeta\) with \(\sigma_*\le\sigma_*'\) and
\(\sigma_*'-\sigma_*\ge(\zeta-\xi-2\eta)_+\); the reverse inclusion bounds
the \(\liminf\) from below by
\[
   \int_{\partial^*\{u>s\}}
      \frac{(\zeta-\xi-2\eta)_+}{|\nabla u|}\,d\Hn .
\]
Letting \(\eta\downarrow0\) gives \eqref{eq:varcoarea}.  The exceptional
levels are fixed as follows: choose the outer and inner simple sandwiches for
the dyadic sequence \(\eta=2^{-m}\), \(m\ge1\), by a fixed dyadic
quantisation of the bounded functions \(\xi,\zeta\).  The union of the
exceptional null sets for those countably many simple pairs is still null and
depends only on \((\xi,\zeta)\).
\end{proof}

\subsection*{Shell-admissible radii}

Let \(\mathcal U\) be the countable family of finite unions of rational balls
whose closures are compactly contained in \(\Om\).  This family is cofinal:
for every compact \(K\Subset\Om\) there is \(U\in\mathcal U\) with
\(K\subset U\Subset\Om\).  For
\[
   c,p,q,\lambda,r\in\mathbb Q,\qquad
   -1\le c\le0,\quad p\le q,\quad
   1\le\lambda\le\rstar^{-1},\quad r>0,\qquad z,a\in\mathbb Q^n,
\]
and \(U\in\mathcal U\), register the bounded Borel pairs
\[
   (p{\bf1}_U,q{\bf1}_U)
\]
and
\[
   \left(
      {\bf1}_U\bigl(\min(c,\nabla u\cdot(\lambda(x-z)+a))-r\bigr),\,
      {\bf1}_U\bigl(\max(c,\nabla u\cdot(\lambda(x-z)+a))+r\bigr)
   \right).
\]
This is a fixed countable family, independent of the radius being tested.
We call \(\rho\in G_\tau\cap(0,1)\) \emph{shell-admissible} if:
\begin{enumerate}[label=(\roman*),leftmargin=2em]
\item \(E_\rho\) has finite perimeter,
\(\partial^*E_\rho\subset\Om\) up to \(\Hn\)-null sets, and
\(\Hn(\partial^*E_\rho\cap\{|\nabla u|=0\})=0\);
\item \(t\) is differentiable at \(\rho\), and
\[
   \wt(\rho)\int_{\partial^*E_\rho}\frac1{|\nabla u|}\,d\Hn
   =\Per(B_\rho);
\]
\item \(s=t(\rho)\) is a level where Lemma~\ref{lem:varcoarea} holds for
every registered pair above.
\end{enumerate}
Almost every \(\rho\in G_\tau\cap(0,1)\) is shell-admissible.  Indeed, each
registered pair has an exceptional null set of levels \(s\).  If \(N\) is
such a null set, then the one-dimensional area formula for the monotone
absolutely continuous map \(t\) gives, on \(G_\tau\),
\[
   |\{\rho:t(\rho)\in N\}|
   \le \tau_0^{-1}\int_{\{t(\rho)\in N\}} \wt(\rho)\,d\rho
   \le \tau_0^{-1}|N|=0.
\]
The remaining conditions hold for a.e.\ radius by BV coarea and the
finite-perimeter level-set package.  The critical-boundary condition follows
from the preceding interior critical-set nullity and the coarea formula:
\[
   \int_0^{\|u\|_\infty}
      \Hn(\partial^*\{u>s\}\cap\{|\nabla u|=0\})\,ds
   =
   \int_{\{|\nabla u|=0\}}|\nabla u|\,dx=0 .
\]

\begin{lemma}[Local BV deformation tail]\label{lem:BV-tail}
Let \(E\subset B_R\) be a set of finite perimeter, let
\(\eta\in C_c^1(\R^n)\) satisfy \(0\le\eta\le1\), and let \(W\in C^1\) on a
neighbourhood of a ball containing \(E\cup\operatorname{spt}\eta\).  Set
\(T_h(x)=x+hW(x)\).  Then
\[
   \limsup_{h\to0}\frac1{|h|}
   \int_{\R^n}\eta(x)
      \bigl|{\bf1}_{T_h(E)}(x)-{\bf1}_E(x)\bigr|\,dx
   \le
   \int_{\partial^*E}\eta\,|W|\,d\Hn .
\]
\end{lemma}

\begin{proof}
For \(|h|\) small, \(T_h\) is a \(C^1\) diffeomorphism on the fixed ball.
Changing variables \(x=T_hy\), it is enough, up to a multiplicative
\(1+o(1)\) factor, to estimate
\[
   \int \eta(T_hy)\bigl|{\bf1}_E(T_hy)-{\bf1}_E(y)\bigr|\,dy .
\]
Let \(u_\eps={\bf1}_E*\varphi_\eps\).  For fixed \(h\), the corresponding
integrals with \(u_\eps\) converge to the displayed integral because
\(u_\eps\to{\bf1}_E\) in \(L^1\) and composition with \(T_h\) has bounded
Jacobian.  For smooth \(u_\eps\),
\[
   |u_\eps(T_hy)-u_\eps(y)|
   \le |h|\int_0^1
      |\nabla u_\eps(y+\theta hW(y))|\,|W(y)|\,d\theta .
\]
After integrating, change variables \(x=y+\theta hW(y)\).  These maps have
Jacobians \(1+o(1)\) uniformly in \(\theta\).  Thus, for fixed \(h\),
\[
   \frac1{|h|}\int \eta(T_hy)|u_\eps(T_hy)-u_\eps(y)|\,dy
   \le
   (1+o_h(1))\int \psi_h\,d|Du_\eps|,
\]
where the continuous compactly supported weights \(\psi_h\) are uniformly
bounded and satisfy \(\psi_h\to\eta|W|\) uniformly as \(h\to0\).  Letting
\(\eps\downarrow0\) at fixed \(h\), strict \(BV_{\rm loc}\) convergence of
the standard mollifications and Reshetnyak continuity for continuous weights
\cite{AFP2000,Maggi2012} give the same bound with \(d|D{\bf1}_E|\).  Taking
\(\limsup_{h\to0}\) then gives
\[
   \limsup_{h\to0}\frac1{|h|}
   \int \eta(x)\bigl|{\bf1}_{T_h(E)}(x)-{\bf1}_E(x)\bigr|\,dx
   \le
   \int \eta|W|\,d|D{\bf1}_E|.
\]
Since \(|D{\bf1}_E|=\Hn\llcorner\partial^*E\), this is the claim.
\end{proof}

\section{Affine moving-ball first variation}\label{sec:shell}

\begin{lemma}[Affine shell estimate]\label{lem:shell}
Let \(\rho\) be shell-admissible.  Fix \(z,a\in\R^n\), and define
\[
   T_h(x):=z+ha+\left(1+\frac h\rho\right)(x-z)
   =x+h\left(\frac{x-z}{\rho}+a\right).
\]
Then
\begin{equation}\label{eq:shell}
   \limsup_{h\downarrow0}
   \frac{|E_{\rho+h}\Delta T_h(E_\rho)|}{h}
   \le
   \int_{\partial^*E_\rho}
      |V_\rho-H_{z,\rho}-a\cdot\nu_\rho|\,d\Hn .
\end{equation}
\end{lemma}

\begin{proof}
Let \(E=E_\rho\), \(W(x)=(x-z)/\rho+a\), and \(w=\wt(\rho)\).  Since
\(\Om\subset B_R\), \(W\) is bounded on \(\Om\).  The finite Radon measure
\[
   \left(
      \frac{w}{|\nabla u|}+|W|
      +\left|\frac{w}{|\nabla u|}-W\cdot\nu_\rho\right|
   \right)\Hn\llcorner\partial^*E
\]
is finite.  Fix \(\eps>0\).  By inner regularity and
\(\partial^*E\subset\Om\) up to \(\Hn\)-null sets, choose a compact
\(K\subset\partial^*E\cap\Om\) such that this measure of
\(\partial^*E\setminus K\) is at most \(\eps\).  Then choose
\(U\in\mathcal U\) with \(K\subset U\Subset\Om\).

We first estimate the symmetric difference in \(U\).  On a neighbourhood of
\(\overline U\), interior elliptic regularity gives \(u\in C^1\), and
uniformly for \(x\in U\),
\[
   T_h^{-1}(x)=x-hW(x)+O(h^2),\qquad
   u(T_h^{-1}x)=u(x)-h\nabla u(x)\cdot W(x)+o(h).
\]
Also \(t(\rho+h)=t(\rho)-hw+o(h)\).  Set
\[
   F(x):=w+\nabla u(x)\cdot W(x).
\]
For \(x\in U\),
\[
   x\in E_{\rho+h}
   \quad\Longleftrightarrow\quad u(x)>t(\rho)-hw+o(h),
\]
while
\[
   x\in T_h(E)
   \quad\Longleftrightarrow\quad u(x)>t(\rho)+h\nabla u(x)\cdot W(x)+o(h).
\]
Thus the part of the symmetric difference in \(U\) is contained, up to a
uniform \(o(h)\) outward error, in the two-sided slab whose endpoints are
\(\min(-w,\nabla u\cdot W)\) and \(\max(-w,\nabla u\cdot W)\).
Precisely, fix \(\delta>0\).  Choose rational
\(c,\lambda,z_0,a_0\) with \(c\in[-1,0]\), \(1\le\lambda\le\rstar^{-1}\),
and
\[
   |c+w|+
   \sup_{\overline U}
   |\nabla u\cdot(\lambda(x-z_0)+a_0)-\nabla u\cdot W|
   \le\delta .
\]
Then choose a rational buffer \(r>2\delta\).  For all sufficiently small
\(h>0\), the local symmetric difference is contained in the registered slab
with endpoints
\[
   \min(c,\nabla u\cdot(\lambda(x-z_0)+a_0))-r,\qquad
   \max(c,\nabla u\cdot(\lambda(x-z_0)+a_0))+r .
\]
Lemma~\ref{lem:varcoarea}, followed by \(\delta,r\downarrow0\), gives
\[
   \limsup_{h\downarrow0}\frac{|(E_{\rho+h}\Delta T_h(E))\cap U|}{h}
   \le
   \int_{\partial^*E\cap U}\frac{|F|}{|\nabla u|}\,d\Hn .
\]

On the complement use
\[
   E_{\rho+h}\Delta T_h(E)\subset (E_{\rho+h}\Delta E)\cup(E\Delta T_h(E)).
\]
For the level-set part, the exact volume derivative and the registered
localized one-sided coarea formula on \(U\) give the complement estimate.
Indeed,
\[
   |E_{\rho+h}\setminus E|
   =\omega_n\bigl((\rho+h)^n-\rho^n\bigr)
   =h\,\Per(B_\rho)+o(h)
   =
   h\,w\int_{\partial^*E}\frac1{|\nabla u|}\,d\Hn+o(h),
\]
where the last equality is the shell-admissible flux identity.  On \(U\),
the same rational-buffer use of Lemma~\ref{lem:varcoarea}, now with one
endpoint \(0\), gives
\[
   \lim_{h\downarrow0}
   \frac{|(E_{\rho+h}\setminus E)\cap U|}{h}
   =
   w\int_{\partial^*E\cap U}\frac1{|\nabla u|}\,d\Hn .
\]
Subtracting the local shell from the exact total shell yields
\[
   \limsup_{h\downarrow0}
   \frac{|(E_{\rho+h}\Delta E)\setminus U|}{h}
   \le
   w\int_{\partial^*E\setminus U}\frac1{|\nabla u|}\,d\Hn
   \le
   w\int_{\partial^*E\setminus K}\frac1{|\nabla u|}\,d\Hn .
\]
For the affine deformation part, choose \(\eta\in C_c^1(\R^n)\) with
\(\eta=0\) on a neighbourhood \(N_K\Subset U\) of \(K\), \(\eta=1\) on a fixed
neighbourhood of \(\overline{B_{R+1}}\setminus U\), and \(0\le\eta\le1\).
Since \(E\subset B_R\) and \(W\) is bounded, \(T_h(E)\subset B_{R+1}\) for all
small \(h>0\); hence \((E\Delta T_h(E))\setminus U\subset
\overline{B_{R+1}}\setminus U\subset\{\eta=1\}\).  As \(\eta=0\) near \(K\) and
\(0\le\eta\le1\), Lemma~\ref{lem:BV-tail} gives
\[
   \limsup_{h\downarrow0}
   \frac{|(E\Delta T_h(E))\setminus U|}{h}
   \le
   \int_{\partial^*E}\eta\,|W|\,d\Hn
   \le
   \int_{\partial^*E\setminus K}|W|\,d\Hn .
\]
The two tail terms are bounded by \(2\eps\).  Therefore
\[
   \limsup_{h\downarrow0}\frac{|E_{\rho+h}\Delta T_h(E_\rho)|}{h}
   \le
   \int_{\partial^*E_\rho\cap U}\frac{|F|}{|\nabla u|}\,d\Hn+2\eps
   \le
   \int_{\partial^*E_\rho}\frac{|F|}{|\nabla u|}\,d\Hn+2\eps .
\]
Letting \(\eps\downarrow0\), and using
\[
   \frac{F}{|\nabla u|}
   =
   \frac{w}{|\nabla u|}
   +\frac{\nabla u\cdot W}{|\nabla u|}
   =
   V_\rho-W\cdot\nu_\rho
   =
   V_\rho-H_{z,\rho}-a\cdot\nu_\rho
\]
on \(\partial^*E_\rho\), proves the estimate.
\end{proof}

\begin{theorem}[One-sided first variation of \(q\)]\label{thm:qplus}
For every \(\rho\in(0,1]\), define
\[
   q(\rho):=\inf_{z\in\R^n}|E_\rho\Delta B_\rho(z)|.
\]
At every shell-admissible \(\rho\in(0,1)\),
\begin{equation}\label{eq:qplus}
   D^+q(\rho):=
   \limsup_{h\downarrow0}\frac{q(\rho+h)-q(\rho)}{h}
   \le
   \frac{n}{\rho}q(\rho)
   +
   \inf_{a\in\R^n}\int_{\partial^*E_\rho}
      |V_\rho-H_{z_\rho,\rho}-a\cdot\nu_\rho|\,d\Hn ,
\end{equation}
where \(z_\rho\) is any minimiser of
\(z\mapsto |E_\rho\Delta B_\rho(z)|\).
\end{theorem}

\begin{proof}
The minimiser exists.  Indeed, \(z\mapsto |E_\rho\Delta B_\rho(z)|\) is
continuous by translation continuity in \(L^1\), and if
\(|z|>R+\rho\), then \(B_\rho(z)\cap E_\rho=\emptyset\), so the value is
\(2|E_\rho|\).  Thus the infimum is attained in a closed ball.

Fix a minimiser \(z_\rho\) and \(a\in\R^n\).  For the affine map in
Lemma~\ref{lem:shell},
\[
   T_h(B_\rho(z_\rho))=B_{\rho+h}(z_\rho+ha).
\]
Therefore
\[
\begin{aligned}
   q(\rho+h)
   &\le |E_{\rho+h}\Delta T_h(B_\rho(z_\rho))|\\
   &\le |E_{\rho+h}\Delta T_h(E_\rho)|
      +|T_h(E_\rho)\Delta T_h(B_\rho(z_\rho))|.
\end{aligned}
\]
Since \(T_h\) is injective affine with constant Jacobian
\((1+h/\rho)^n\),
\[
   |T_h(E_\rho)\Delta T_h(B_\rho(z_\rho))|
   =
   \left(1+\frac h\rho\right)^n q(\rho).
\]
Subtract \(q(\rho)\), divide by \(h\), take the limsup as \(h\downarrow0\),
and use Lemma~\ref{lem:shell}.  Then take the infimum over \(a\).
\end{proof}

\begin{lemma}[Unconditional Lipschitz bound for \(q\)]\label{lem:qLip}
For \(0<\rho\le\sigma\le1\),
\[
   |q(\sigma)-q(\rho)|
   \le 2\omega_n(\sigma^n-\rho^n)
   \le 2n\omega_n|\sigma-\rho|.
\]
Consequently \(q\) is absolutely continuous and differentiable a.e.
\end{lemma}

\begin{proof}
Use the same centre for the two radii and the nestedness
\(E_\rho\subset E_\sigma\), \(B_\rho(z)\subset B_\sigma(z)\):
\[
   q(\sigma)\le q(\rho)+|E_\sigma\setminus E_\rho|
                         +|B_\sigma\setminus B_\rho|
   =q(\rho)+2\omega_n(\sigma^n-\rho^n).
\]
Interchange \(\rho\) and \(\sigma\) to get the reverse inequality.
\end{proof}

\section{Positive kinetic estimate}\label{sec:kinetic}

\begin{lemma}[Transfer to a distance-minimising centre]\label{lem:center-transfer}
After reducing \(\eta_I=\eta_I(n,R,\rstar)\) if necessary, the following
holds.  Let \(\rho\in G_I\) be a level where
Theorem~\ref{inp:samecenter} holds, and let \(z_\rho^{\min}\) minimise
\(z\mapsto |E_\rho\Delta B_\rho(z)|\).  Then
\begin{equation}\label{eq:H-DI-min}
   \left(
   \int_{\partial^*E_\rho}
      |H_{z_\rho^{\min},\rho}-\bar V_\rho|\,d\Hn
   \right)^2
   \le C(n,R,\rstar)D_I(t(\rho)).
\end{equation}
\end{lemma}

\begin{proof}
Let \(z_\rho^{\rm sc}\) be the centre from
Theorem~\ref{inp:samecenter}.  Minimality gives
\[
   |E_\rho\Delta B_\rho(z_\rho^{\min})|
   \le |E_\rho\Delta B_\rho(z_\rho^{\rm sc})|
   \le C D_I(t(\rho))^{1/2}.
\]
Thus
\[
   |B_\rho(z_\rho^{\min})\Delta B_\rho(z_\rho^{\rm sc})|
   \le C D_I(t(\rho))^{1/2}.
\]
For equal balls with centre distance \(d\),
\[
   |B_\rho(z_1)\Delta B_\rho(z_2)|
   \ge c(n)\rho^{n-1}\min\{d,\rho\}.
\]
Indeed, in coordinates \(x=z_1+s\,e+y\) with \(e=(z_2-z_1)/d\) and
\(y\perp e\), one has \(B_\rho(z_1)\setminus B_\rho(z_2)
=\{s^2+|y|^2<\rho^2,\ s\le d/2\}\); integrating the \((n-1)\)-dimensional
cross-sections over \(s\in[0,\tfrac12\min\{d,\rho\}]\), where
\(\rho^2-s^2\ge\tfrac34\rho^2\), gives
\(|B_\rho(z_1)\setminus B_\rho(z_2)|\ge
\tfrac12\omega_{n-1}(3/4)^{(n-1)/2}\rho^{n-1}\min\{d,\rho\}\), and the
symmetric difference is twice this.
Choose \(\eta_I\) so small that the preceding upper bound forces the
linear branch \(d<\rho\).  Since \(\rho\ge\rstar\),
\[
   |z_\rho^{\min}-z_\rho^{\rm sc}|
   \le C(n,\rstar)D_I(t(\rho))^{1/2}.
\]
The homothetic trace is affine in the centre:
\[
   H_{z_\rho^{\min},\rho}-H_{z_\rho^{\rm sc},\rho}
   =
   -\frac{(z_\rho^{\min}-z_\rho^{\rm sc})\cdot\nu_\rho}{\rho}.
\]
On \(G_I\), the perimeter is uniformly bounded because
\(D_I(t(\rho))\le\eta_I\) and \(\rho\in[\rstar,1]\).  Hence
\[
   \int_{\partial^*E_\rho}
   |H_{z_\rho^{\min},\rho}-H_{z_\rho^{\rm sc},\rho}|\,d\Hn
   \le C(n,R,\rstar)D_I(t(\rho))^{1/2}.
\]
Combining this with \eqref{eq:sc-input} via the triangle inequality,
\[
   \int_{\partial^*E_\rho}
      |H_{z_\rho^{\min},\rho}-\bar V_\rho|\,d\Hn
   \le
   \int_{\partial^*E_\rho}
      |H_{z_\rho^{\min},\rho}-H_{z_\rho^{\rm sc},\rho}|\,d\Hn
   +\int_{\partial^*E_\rho}
      |H_{z_\rho^{\rm sc},\rho}-\bar V_\rho|\,d\Hn
   \le C(n,R,\rstar)D_I(t(\rho))^{1/2},
\]
and squaring gives \eqref{eq:H-DI-min}.
\end{proof}

\begin{proposition}[Positive kinetic bound]\label{prop:positivekinetic}
Let \(J\subset[\rstar,\rdel]\) be an interval.  Then
\[
   \int_J (q'_+(\rho))^2\,d\rho
   \le C(n,R,\rstar)\,\dT(\Om),
\]
where \(q'_+(\rho)=\max\{q'(\rho),0\}\) at a.e. differentiability points.
\end{proposition}

\begin{proof}
Let
\[
   G:=J\cap G_I\cap G_\tau,\qquad \mathcal B:=J\setminus G .
\]
Since \(q\) is Lipschitz, \(q'\) exists a.e.; shell-admissible radii have
full measure in \(G_\tau\).  At points where both properties hold,
\(D^+q=q'\).  The right side of Theorem~\ref{thm:qplus} is nonnegative,
so the theorem also bounds \(q'_+\).  On \(G\), outside this null set,
Theorem~\ref{thm:qplus}, Theorem~\ref{inp:samecenter},
Lemma~\ref{lem:center-transfer}, and Lemma~\ref{lem:velocity} give
\[
   (q'_+)^2
   \le C(n,R,\rstar)\bigl(D_H(t(\rho))+D_I(t(\rho))\bigr).
\]
Indeed \(q^2\le C D_I\) on \(G_I\), and
\[
\begin{aligned}
   \inf_a\int |V_\rho-H_{z_\rho^{\min},\rho}-a\cdot\nu_\rho|
   &\le
   \int |V_\rho-\bar V_\rho|
   +\int |H_{z_\rho^{\min},\rho}-\bar V_\rho|.
\end{aligned}
\]
On \(G_\tau\), \(d\rho\le\tau_0^{-1}d\mu\).  Hence the budget
\eqref{eq:budget} yields
\[
   \int_G(q'_+)^2\,d\rho\le C\,\dT(\Om).
\]

It remains to pay \(\mathcal B\).  By Lemma~\ref{lem:qLip},
\(q'_+\le2n\omega_n\) a.e.  The bad-density part has
\(|J\cap B_\tau|\le C\,\dT(\Om)\).  For the bad-isoperimetric part on
good density,
\[
   |J\cap G_\tau\cap G_I^c|
   \le \tau_0^{-1}\mu(G_I^c)
   \le \frac{C}{\eta_I}\int_JD_I(t(\rho))\,d\mu(\rho)
   \le C\,\dT(\Om).
\]
Therefore \(|\mathcal B|\le C\,\dT(\Om)\), and
\[
   \int_{\mathcal B}(q'_+)^2\,d\rho\le C\,\dT(\Om).
\]
\end{proof}

\section{Endpoint trace and bounded stability}\label{sec:endpoint}

\begin{lemma}[Moving-ball endpoint trace]\label{lem:endpoint}
Let
\[
   L_0:=\frac{1-\rstar}{8},\qquad
   J:=[\rdel-L_0,\rdel].
\]
Assume \(\dT(\Om)\) is small enough that
\(\rdel\ge(1+\rstar)/2\).  Then \(J\subset[\rstar,\rdel]\) and
\[
   q(\rdel)^2\le C(n,R,\rstar)\,\dT(\Om).
\]
\end{lemma}

\begin{proof}
The inclusion \(J\subset[\rstar,\rdel]\) follows from
\(\rdel-\rstar\ge(1-\rstar)/2\).  On \(G=J\cap G_I\cap G_\tau\),
\[
   q(\rho)^2\le C D_I(t(\rho)),\qquad
   d\rho\le\tau_0^{-1}d\mu(\rho),
\]
so
\[
   \int_G q(\rho)^2\,d\rho\le C\,\dT(\Om).
\]
On \(J\setminus G\), the measure is \(O(\dT)\) by the proof of
Proposition~\ref{prop:positivekinetic}, while \(q\le2\omega_n\).  Hence
\[
   \int_J q(\rho)^2\,d\rho\le C\,\dT(\Om).
\]
There is therefore \(\rho_0\in J\) with
\[
   q(\rho_0)^2\le C L_0^{-1}\dT(\Om).
\]
Since \(q\) is absolutely continuous and \(\rho_0\le\rdel\),
\[
   q(\rdel)
   \le q(\rho_0)+\int_{\rho_0}^{\rdel}q'_+(\rho)\,d\rho.
\]
By Cauchy--Schwarz and Proposition~\ref{prop:positivekinetic},
\[
   q(\rdel)
   \le C\dT(\Om)^{1/2}
      +L_0^{1/2}\left(\int_J(q'_+)^2\,d\rho\right)^{1/2}
   \le C\dT(\Om)^{1/2}.
\]
\end{proof}

\begin{theorem}[Bounded sharp Saint--Venant stability]
\label{thm:boundedSV}
There are computable constants \(C=C(n,R,\rstar)\) and
\(\delta_0=\delta_0(n,R,\rstar)>0\), depending only on the level-set
identity constants, the Fusco--Julin constant, and the explicit constants
above, such that every bounded open \(\Om\subset B_R\),
\(|\Om|=\omega_n\), satisfies
\[
   \Fr(\Om)^2\le C\,\dT(\Om).
\]
\end{theorem}

\begin{proof}
Let \(\kappa\) be the constant fixed in Proposition~\ref{prop:levelset}.
First assume \(\dT(\Om)\le\delta_0\), with \(\delta_0\) chosen so that
\(\rdel=1-\kappa\sqrt{\dT(\Om)}\ge(1+\rstar)/2\), \(\dT(\Om)\le1\), and
the smallness thresholds used in Theorem~\ref{inp:samecenter} and
Lemma~\ref{lem:center-transfer} hold.
Lemma~\ref{lem:endpoint} gives
\[
   q(\rdel)^2\le C\dT(\Om).
\]
Since \(|E_{\rdel}|=\omega_n\rdel^n\) and \(\rdel\ge\rstar\),
\[
   \Fr(E_{\rdel})^2
   \le \frac{q(\rdel)^2}{\omega_n^2\rdel^{2n}}
   \le C(n,\rstar)\dT(\Om).
\]
Moreover \(E_{\rdel}\subset\Om\), and by the definition of \(\rdel\),
\[
   |\Om\setminus E_{\rdel}|
   =\omega_n(1-\rdel^n)
   \le n\omega_n(1-\rdel)
   =n\omega_n\kappa\,\dT(\Om)^{1/2}.
\]
Lemma~\ref{lem:superlevel-transfer} gives
\[
   \Fr(\Om)
   \le \Fr(E_{\rdel})+C(n,\rstar)\dT(\Om)^{1/2}.
\]
Squaring proves the estimate in the small-deficit regime.

If \(\dT(\Om)>\delta_0\), then \(\Fr(\Om)<2\), so
\[
   \Fr(\Om)^2\le \frac{4}{\delta_0}\dT(\Om).
\]
Enlarging \(C\) completes the proof.
\end{proof}

\section{Proof architecture and constants}\label{sec:architecture}

The argument is organised around the following five steps.
\begin{enumerate}[leftmargin=2em]
\item The exact level-set deficit identity (Theorem~\ref{thm:id-main}) writes
\(\dT(\Om)\) as a weighted reduced-boundary integral of the Cauchy--Schwarz
defect \(D_H\) and the isoperimetric defect \(D_I\); the profile-gap identity
and the bad-density bound of Proposition~\ref{prop:levelset} are its
\(\rho\)-parametrised consequences.
\item The affine shell estimate (Lemma~\ref{lem:shell}) is a finite-perimeter
limsup estimate, proved by localized coarea differentiation on variable slabs
(Lemma~\ref{lem:varcoarea}) together with the local BV deformation-tail lemma
(Lemma~\ref{lem:BV-tail}).  Its only bounded-variation input is Reshetnyak
continuity for a continuous weight, applied as an \emph{upper} bound to a
swept-volume tail made arbitrarily small by inner regularity; no
lower-semicontinuity passage is used.
\item The endpoint trace (Lemma~\ref{lem:endpoint}) uses the upper Dini
derivative \(D^+q\) together with the global Lipschitz bound for \(q\)
(Lemma~\ref{lem:qLip}).
\item Bad-density levels are charged by Lebesgue measure rather than by the
deficit, which is possible because \(q\) is uniformly Lipschitz by nestedness.
\item The same-centre good-level package (Theorem~\ref{inp:samecenter}) is
obtained from the Fusco--Julin strong quantitative isoperimetric inequality
(Theorem~\ref{thm:FJ-strong}); the passage from normal oscillation to the
homothetic trace is the radial trace Lemma~\ref{lem:oriented-radial-trace}.
\end{enumerate}

\noindent
All constants in Theorem~\ref{thm:boundedSV} depend only on \(n\), \(R\), and
\(\rstar\), together with the Fusco--Julin constant \(C_{\rm FJ}(n)\).  Tracing
them through the proof: \(C_L=2n^2\omega_n\) and
\(C_\tau=4n/(\omega_n\rstar^{n+1})\) (Proposition~\ref{prop:levelset});
\(\kappa=(1-\rstar)/4\) and \(\tau_0=\rstar/(2n)\); the same-centre constant
\(C_{\rm sc}(n,R,\rstar)\) and threshold \(\eta_I(n,R,\rstar)\)
(Theorem~\ref{inp:samecenter}, Lemma~\ref{lem:center-transfer}); the ball--ball
lower constant \(\tfrac12\omega_{n-1}(3/4)^{(n-1)/2}\); the slab length
\(L_0=(1-\rstar)/8\); and the large-deficit cutoff \(4/\delta_0\).  No constant
depends on \(\Om\) beyond these parameters.

\appendix

\section{Standard tools from the theory of sets of finite perimeter}
\label{app:bv}

For the reader's convenience we record the three classical results used above,
in the exact form needed; all hold for sets of finite perimeter with no graph
or manifold regularity.

\begin{enumerate}[leftmargin=2em]
\item \emph{BV coarea formula} \cite[Thm.~13.1]{Maggi2012},
\cite[Thm.~3.40]{AFP2000}: for \(u\in BV_{\rm loc}(\R^n)\) and Borel
\(\varphi\ge0\),
\(\int|\nabla u|\varphi=\int_{\R}(\int_{\partial^*\{u>t\}}\varphi\,d\Hn)\,dt\).
This underlies \eqref{eq:id-bvcoarea} and Lemma~\ref{lem:varcoarea}.
\item \emph{Gauss--Green on sets of finite perimeter}
\cite[Ch.~17]{Maggi2012}: for a bounded Lipschitz vector field \(X\) and a set
\(E\) of finite perimeter,
\(\int_E\operatorname{div}X=\int_{\partial^*E}X\cdot\nu_E\,d\Hn\).  This is the
form used in Lemma~\ref{lem:oriented-radial-trace}.
\item \emph{Strict \(BV\) convergence and Reshetnyak continuity}
\cite[Thm.~2.39]{AFP2000}, \cite[Ch.~20]{Maggi2012}: if
\(v_k\to v\) in \(L^1_{\rm loc}\) with \(|Dv_k|(\R^n)\to|Dv|(\R^n)\), then
\(\int\psi\,d|Dv_k|\to\int\psi\,d|Dv|\) for every continuous compactly
supported \(\psi\).  Applied to mollifications \(v_k={\bf1}_E*\varphi_{1/k}\),
this gives the upper bound in Lemma~\ref{lem:BV-tail}.
\end{enumerate}

\section{Dependency certificate}\label{app:constants}

Every step is either proved in this manuscript or quoted from a named
published theorem.

\emph{Proved internally:} the exact deficit identity and its inputs (coarea
form of \(-m'\), the Lipschitz-truncation flux identity, the rearranged
primitive with \(I''=-1/\alpha\)) in \S\ref{sec:identity}; the profile-gap and
bad-density estimates in \S\ref{sec:profilegap}; the
oriented radial trace, the same-centre package, the velocity defect, and the
superlevel asymmetry transfer in \S\ref{sec:samecentre}; the two-sided coarea
differentiation lemma, the shell-admissible registration, and the local BV
deformation tail in \S\ref{sec:coarea}; the affine shell estimate, the
one-sided first variation, and the Lipschitz bound for \(q\) in
\S\ref{sec:shell}; the centre transfer and the positive kinetic estimate in
\S\ref{sec:kinetic}; the endpoint trace and the bounded stability theorem in
\S\ref{sec:endpoint}.

\emph{Quoted with exact citation:} the Fusco--Julin strong quantitative
isoperimetric inequality \cite{FJ2014}; the BV coarea formula, the
finite-perimeter Gauss--Green theorem, and Reshetnyak continuity
\cite{AFP2000,Maggi2012}; the De~Giorgi isoperimetric inequality
\cite{DeGiorgi1958}; the Talenti rearrangement \(L^\infty\) bound
\cite{Talenti1976}; the absolutely continuous monotone inverse theorem
\cite{BBC1975}; Stampacchia's theorem on \(\{\nabla u=0\}\)
\cite{Stampacchia1965}; interior elliptic regularity and the weak maximum
principle \cite{GilbargTrudinger2001}; and the P\'olya--Szeg\H{o}
identification of the ball energy \cite{PolyaSzego1951,Bandle1980}.  Every
ingredient is therefore either proved above or quoted from the published
literature; the argument can be checked from this manuscript alone.

\begin{thebibliography}{99}

\bibitem{AFP2000}
L.~Ambrosio, N.~Fusco, and D.~Pallara,
\emph{Functions of Bounded Variation and Free Discontinuity Problems},
Oxford University Press, 2000.
(Thm.~3.40, coarea; Thm.~3.99, chain rule; \S3.5, reduced boundaries.)

\bibitem{Bandle1980}
C.~Bandle,
\emph{Isoperimetric Inequalities and Applications},
Pitman, 1980. Ch.~II.

\bibitem{BBC1975}
P.~B\'enilan, H.~Brezis, and M.~G.~Crandall,
\emph{A semilinear equation in \(L^1(\R^N)\)},
J.~Anal.\ Math.\ \textbf{27} (1975), 273--299.
(Thm.~0.4 and \S1: absolutely continuous monotone inverse.)

\bibitem{DeGiorgi1958}
E.~De~Giorgi,
\emph{Sulla propriet\`a isoperimetrica dell'ipersfera, nella classe degli
insiemi aventi frontiera orientata di misura finita},
Atti Accad.\ Naz.\ Lincei \textbf{5} (1958), 33--44.

\bibitem{FJ2014}
N.~Fusco and V.~Julin,
\emph{A strong form of the quantitative isoperimetric inequality},
Calc. Var. Partial Differential Equations \textbf{50} (2014), 925--937.

\bibitem{GilbargTrudinger2001}
D.~Gilbarg and N.~S.~Trudinger,
\emph{Elliptic Partial Differential Equations of Second Order},
Classics in Mathematics, Springer, 2001.
(Interior \(W^{2,p}\) regularity; weak maximum principle.)

\bibitem{Maggi2012}
F.~Maggi,
\emph{Sets of Finite Perimeter and Geometric Variational Problems},
Cambridge University Press, 2012.
(Ch.~13, coarea; Ch.~17, Gauss--Green on finite-perimeter sets.)

\bibitem{PolyaSzego1951}
G.~P\'olya and G.~Szeg\H{o},
\emph{Isoperimetric Inequalities in Mathematical Physics},
Princeton University Press, 1951.

\bibitem{Stampacchia1965}
G.~Stampacchia,
\emph{Le probl\`eme de Dirichlet pour les \'equations elliptiques du second
ordre \`a coefficients discontinus},
Ann.\ Inst.\ Fourier \textbf{15} (1965), 189--258.
(\(D^2u=0\) a.e.\ on \(\{\nabla u=0\}\) for Sobolev maps.)

\bibitem{Talenti1976}
G.~Talenti,
\emph{Elliptic equations and rearrangements},
Ann.\ Mat.\ Pura Appl.\ \textbf{110} (1976), 353--372.

\end{thebibliography}

\end{document}
